\chapter{Le regole che insegnano alla rete: la distillazione}\label{cap:p1-distillazione}

% Capitolo della Parte I --- la knowledge distillation spiegata in modo semplice.

Abbiamo due protagonisti: l'\textbf{esperto} che conosce le regole (il teacher)
e la \textbf{scatola che impara} (la rete neurale). La rete non ha idee
proprie: ha semplicemente copiato le risposte che l'esperto dava sulle
flashcard. Questo passaggio si chiama \textbf{knowledge distillation}
(distillazione della conoscenza): la conoscenza dell'esperto viene ``distillata''
e trasferita dentro la rete.

\section{Perché non ci si ferma alle regole?}\label{sec:p1-dist-perche}

Se l'esperto sa già rispondere, perché costruire una rete che lo imita? Perché
le due hanno vantaggi diversi:

\begin{itemize}
  \item \textbf{Le regole} sono chiare e sicure, ma vanno scritte a mano e sono
        rigide: non gestiscono bene i casi nuovi o sfumati.
  \item \textbf{La rete} è veloce, automatica e impara dai dati, ma può
        sbagliare e non spiega il perché di una scelta.
\end{itemize}

La distillazione prende il meglio di entrambe: la rete impara la \emph{logica}
dell'esperto (così diventa flessibile e veloce), ma le regole restano la
\emph{fonte} da cui tutto deriva. Nel progetto le due convivono: la rete fa il
lavoro di tutti i giorni, le regole restano il punto di riferimento.

\begin{esempio}{Nel progetto}
La rete (MLP) viene addestrata sulle 70\,000 flashcard etichettate dall'esperto.
Alla fine, se le mostri un paziente nuovo, risponde con il rischio e con un
numero di ``sicurezza'' (quanto si fida della sua risposta). Le regole, però,
restano lì pronte a intervenire quando serve --- ed è il tema del prossimo
capitolo.
\end{esempio}
